\section{Delimitación}
Dentro de ese tema, el proyecto se delimita a la validación de identidad en
un acceso peatonal del Tecnológico de Estudios Superiores de Chalco (TESCHA),
en Chalco, Estado de México. Se elige un solo acceso, en el turno vespertino,
como punto piloto tanto para la observación inicial como para la prueba del
prototipo (Apéndices~A, B y~C), y se opta por el código QR dinámico como
tecnología de identificación, según se justifica en ``La Viabilidad'', más
adelante.

Concretamente, se diseñará y construirá un prototipo de torniquete
controlado por un microcontrolador ESP32 con cámara (ESP32-CAM), capaz de
leer un código QR dinámico asignado a cada estudiante, validarlo contra un
servidor y registrar cada ingreso en una base de datos consultable desde un
panel web (Figura~\ref{fig:arquitectura}). El código se regenera cada pocos
segundos y se firma con una contraseña de un solo uso basada en HMAC, un
esquema que resiste la interceptación, el reenvío y la modificación del
mensaje \parencite{subpratatsavee2015hmac}; así, una captura de pantalla
compartida deja de ser válida al poco tiempo. Esta estrategia de regeneración
periódica ya se ha probado en el control de asistencia de una universidad del
Reino Unido, donde un código QR renovado cada 20 segundos impidió que los
estudiantes reenviaran capturas de pantalla entre sí
\parencite{nwabuwe2023dynamicqr}; otros trabajos han combinado de forma
similar el QR con una contraseña de un solo uso para el acceso a laboratorios
universitarios \parencite{chaabani2023labqr}.

\begin{figure}[tbp]
	\caption{Arquitectura Propuesta del Sistema de Acceso por Código QR}
	\label{fig:arquitectura}
	\centering
	\begin{tikzpicture}[
		node distance=1.1cm and 1.6cm,
		box/.style={draw, rounded corners, align=center, minimum height=1cm,
				minimum width=2.6cm, font=\small},
		arr/.style={-{Latex[length=2mm]}, thick}
		]
		\node[box] (alumno) {QR dinámico\\(celular del alumno)};
		\node[box, right=of alumno] (cam) {ESP32-CAM\\(captura y decodifica)};
		\node[box, right=of cam] (esp) {ESP32\\(controla torniquete)};
		\node[box, below=of esp] (servidor) {Servidor y BD\\(valida el token)};
		\node[box, left=of servidor] (panel) {Panel web\\administrativo};
		\draw[arr] (alumno) -- (cam) node[midway, above, font=\scriptsize]{escaneo};
		\draw[arr] (cam) -- (esp) node[midway, above, font=\scriptsize]{token};
		\draw[arr] (esp) -- (servidor) node[midway, right, font=\scriptsize]{Wi-Fi/HTTPS};
		\draw[arr] (servidor) -- (panel) node[midway, below, font=\scriptsize]{consulta};
	\end{tikzpicture}
	\figurenote{Elaboración propia. El servidor responde si el token es válido
		y el ESP32 libera el torniquete; cada intento, aceptado o rechazado, queda
		registrado en la base de datos.}
\end{figure}

La delimitación del proyecto en esta etapa es la siguiente:
\begin{itemize}
	\item \textbf{Incluye:} un acceso piloto, el registro de entradas de
	      estudiantes, la generación y validación de QR dinámicos, la base de datos de
	      ingresos y un panel web de consulta.
	\item \textbf{No incluye:} reconocimiento facial, tarjetas NFC, pase de lista
	      por materia, registro de salidas ni control de visitantes. Los estudiantes
	      que no puedan mostrar su código seguirán el procedimiento manual actual.
\end{itemize}

Cabe reconocer una limitación: un QR dinámico impide reutilizar capturas de
pantalla, pero no que un estudiante preste su celular a otra persona. Por esa
razón, el proyecto no pretende eliminar la suplantación, sino ayudar a
aliviar los estragos que genera el sistema actual.
